\subsection{RESTRICTION AND EXTENSION OF SCALARS}

Let \(A_{0}\) and A be two commutative rings and \(\rho:A_{0}\to A\) a ring homomorphism. If E is an A-algebra, we denote (conforming with II, § 1, no. 13) by \(\rho_{*}(E)\) the \(A_{0}\) -module defined by addition on E and the external law

\[
\lambda . x = \rho (\lambda) x \quad \text { for   all } \lambda \in A _ {0} \text { and   all } x \in E.
\]

Multiplication in E and the \(A_{0}\) -module structure on \(\rho_{*}(E)\) define an \(A_{0}\) -algebra structure on \(\rho_{*}(E)\) . When \(A_{0}\) is a subring of A and \(\rho\) the canonical injection, the algebra \(\rho_{*}(E)\) is said to be obtained from E by restricting the ring A of scalars to \(A_{0}\) . By an abuse of language, this is also sometimes said when \(\rho\) is arbitrary.

Let F be an \(A_{0}\) -algebra. An \(A_{0}\) -algebra homomorphism \(F \to \rho_{*}(E)\) is called a semi-homomorphism (relative to \(\rho\) ) or a \(\rho\) -homomorphism of F into the A-algebra E; it is also called an \(A_{0}\) -homomorphism if no confusion arises. If E, \(E'\) are two A-algebras, every A-algebra homomorphism \(E \to E'\) is also an \(A_{0}\) -algebra homomorphism \(\rho_{*}(E) \to \rho_{*}(E')\) .

Consider now two commutative rings A and B and a ring homomorphism \(\rho: \mathbf{A} \to \mathbf{B}\) . For every A-module E, the B-module \(\rho^{*}(\mathbf{E}) = \mathbf{E} \otimes_{\mathbf{A}} \mathbf{B}\) , obtained from E by extending the ring A of scalars to B, has been defined (II, § 5, no. 1). If E is also an A-algebra, we shall define on \(\rho^{*}(\mathbf{E})\) a B-algebra structure. For this, observe that \((\mathbf{E} \otimes_{\mathbf{A}} \mathbf{B}) \otimes_{\mathbf{B}} (\mathbf{E} \otimes_{\mathbf{A}} \mathbf{B})\) is canonically isomorphic to \((\mathbf{E} \otimes_{\mathbf{A}} \mathbf{E}) \otimes_{\mathbf{A}} \mathbf{B}\) (II, § 5, no. 1, Proposition 3). If \(m: \mathbf{E} \otimes_{\mathbf{A}} \mathbf{E} \to \mathbf{E}\) defines the multiplication on E, the mapping \(m \otimes l_{\mathbf{B}}: (\mathbf{E} \otimes_{\mathbf{A}} \mathbf{E}) \otimes_{\mathbf{A}} \mathbf{B} \to \mathbf{E} \otimes_{\mathbf{A}} \mathbf{B}\) is therefore canonically identified with a B-linear mapping

\[
m ^ {\prime}: \rho^ {*} (\mathbf {E}) \otimes_ {\mathbf {B}} \rho^ {*} (\mathbf {E}) \rightarrow \rho^ {*} (\mathbf {E})
\]

which defines the desired B-algebra structure on \(\rho^{*}(\mathbf{E})\) . Hence

\[
(x \otimes \beta) (x ^ {\prime} \otimes \beta^ {\prime}) = (x x ^ {\prime}) \otimes (\beta \beta^ {\prime})\tag{5}
\]

for \(x, x'\) in \(\mathbf{E}\) , \(\beta\) and \(\beta'\) in \(\mathbf{B}\) . The B-algebra \(\rho^*(\mathbf{E})\) is said to be derived from the A-algebra \(\mathbf{E}\) by extending the ring \(\mathbf{A}\) of scalars to \(\mathbf{B}\) (by means of \(\rho\) ). It is also denoted by \(\mathrm{E}_{(\mathbf{B})}\) or \(\mathbf{E} \otimes_{\mathbf{A}} \mathbf{B}\) . When \(\mathbf{E}\) is associative (resp. commutative, resp. unital), so is \(\rho^*(\mathbf{E})\) .

PROPOSITION 1. For every A-algebra E, the canonical mapping \(\phi_{\mathbf{E}}: x \mapsto x \otimes 1\) of E into \(\mathrm{E}_{(\mathbf{B})}\) is an A-homomorphism of algebras. Moreover, for every B-algebra F and every A-homomorphism \(f: \mathbf{E} \to \mathbf{F}\) , there exists one and only one B-homomorphism \(\bar{f}: \mathbf{E}_{(\mathbf{B})} \to \mathbf{F}\) such that \(\bar{f}(x \otimes 1) = f(x)\) for all \(x \in \mathbf{E}\) .

The first assertion follows immediately from the definition of multiplication in \(\mathrm{E}_{(\mathrm{B})}\) , which gives \((x\otimes1)(x'\otimes1)=(xx')\otimes1\) for \(x\in E\) and \(x'\in E\) . The existence and uniqueness of the B-linear mapping \(\bar{f}\) of \(\mathrm{E}_{(\mathrm{B})}\) into F satisfying the relation \(\bar{f}(x\otimes1)=f(x)\) for all \(x\in E\) follow from II, §5, no. 1, Proposition 1; here it all amounts to verifying that \(\bar{f}(yy')=\bar{f}(y)\bar{f}(y')\) for y and \(y'\) in \(\mathrm{E}_{(\mathrm{B})}\) ; as the elements of the form \(x\otimes1\) (with \(x\in E\) ) generate the B-module \(\mathrm{E}_{(\mathrm{B})}\) , attention may be confined to the case where \(y=x\otimes1\) , \(y'=x'\otimes1\) with \(x\in E\) , \(x'\in E\) ; as \(yy'=(xx')\otimes1\) , the relation \(\bar{f}(yy')=\bar{f}(y)\bar{f}(y')\) then follows from \(f(xx')=f(x)f(x')\) .

It can also be said that \(f \mapsto \bar{f}\) is a canonical bijection

\[
\operatorname{Hom} _ {\mathrm{A-alg.}} (\mathrm{E}, \rho_ {*} (\mathrm{F})) \rightarrow \operatorname{Hom} _ {\mathrm{B-alg.}} (\rho^ {*} (\mathrm{E}), \mathrm{F}).\tag{6}
\]

The ordered pair consisting of \(E_{(B)}\) and \(\phi_{E}\) is therefore a solution of the universal mapping problem (Set Theory, IV, §3, no. 1) where \(\Sigma\) is the species of B-algebra structure and the \(\alpha\) -mappings the A-homomorphisms from E to a B-algebra.

COROLLARY. Let E, E' be two A-algebras; for every A-homomorphism of algebras \(u: E \to E'\) , \(u \otimes l_{B}\) is the unique B-homomorphism of algebras v: \(E \otimes_{A} B \to E' \otimes_{A} B\) rendering commutative the diagram

\[
\begin{array}{c} \mathrm{E} \xrightarrow {\phi_ {\mathrm{E}}} \mathrm{E} \otimes_ {\mathrm{A}} \mathrm{B} \\ u \Biggl \downarrow \\ \mathrm{E} ^ {\prime} \xrightarrow [ \phi_ {\mathrm{E} ^ {\prime}} ]{} \mathrm{E} ^ {\prime} \otimes_ {\mathrm{A}} \mathrm{B} \end{array}
\]

Let C be a third commutative ring and \(\sigma: B \to C\) a ring homomorphism; it is immediate that the canonical C-homomorphism

\[
\sigma^ {*} (\rho^ {*} (\mathbf {E})) \rightarrow (\sigma \circ \rho) ^ {*} (\mathbf {E})
\]

mapping \((x\otimes 1)\otimes 1\) to \(x\otimes 1\) for all \(x\in \mathbf{E}\) (II, §5, no. 1, Proposition 2) is an algebra isomorphism.

